\documentclass[11pt]{ctexart}
\usepackage{fancyhdr}

% 清除页眉内容并去掉页眉横线
\fancyhf{} % 清空所有页眉页脚
\renewcommand{\headrulewidth}{0pt} % 去掉页眉横线
\pagestyle{fancy} % 应用设置
\usepackage{graphicx}
\usepackage{fix-cm}
\usepackage{setspace}
\usepackage{amsmath}
\usepackage{geometry}
\linespread{1.5}
\geometry{left=2.5cm,right=2.5cm,top=2.5cm,bottom=2.5cm}
\title{电磁学笔记}
\author{CCY}
\date{\today}
\begin{document}
\tableofcontents%生成目录
\newpage%新建页，内容和目录分开
\maketitle
\section{基本的电磁现象}
\subsection{电荷与电场}
\subsection{导体与电介质}
\subsection{电流场的描述}
\subsection{磁场}
\subsection{磁矩}
\newpage
\section{电场}
\subsection{高斯定理}
\subsection{电场的散度}
\subsection{静电场的电势}
\subsection{静电势能和电场中的能量}
\subsubsection{电荷在外电场的静电势能}
\subsubsection{带点体系的静电能}
电势能存储在电场中:类比弹簧：势能储存在弹簧中，压缩弹簧，能量
被存储到弹簧中；两个正电荷相互靠近，电场发生改变，电势能存储
到电场中；      不同之处：弹簧的弹性势能与两端的质量块无关，
电势能与电场强度有关

静电势能<-->电势<-->做功<-->电场力：
体系的电势能：$E=\frac{1}{2} \int \rho(\vec{r})\phi(\vec{r})dV$
=$\frac{1}{2} \sum q_iU_i$=$\frac{1}{2} \iint \rho_e u \, d\tau $

\noindent 例题：
\begin{enumerate}
    \item 均匀带电球壳带电量为Q,求它的电势能
        
    解：在空间中，将电荷不断从无穷远处移动到球壳上，使球壳的带电量从0增加至Q,在这个过程中，外力做的功即为体系的电势能：

    运用公式求解：$U=\frac{1}{4 \pi \varepsilon_0} \frac{Q}{R}$
    
    $W=\frac{1}{2}Q U =\frac{1}{8 \pi \varepsilon_0} \frac{Q^2}{R}$

    \item 均匀带电球体，半径为$R$，电荷密度为$\rho$，求它的电势能
    
    解：利用公式求解：

    $\vec{E}=\begin{cases}
        \frac{1}{4 \pi \varepsilon_0} \frac{Q}{r^2} & r<R \\
        \frac{1}{4 \pi \varepsilon_0} \frac{Q}{R^3}  r & r>R
    \end{cases}$
   \vspace{10pt}
   \begin{spacing}{2.0}

    $U=\frac{1}{4 \pi \varepsilon_0} (\int_r^R \vec{E} dx+ \int_R ^\infty \vec{E} dx)
    =\frac{1}{4 \pi \varepsilon_0 R^3} \int_r^R x dx+ \int_R ^\infty \frac{R^3}{x^2} dx
    =\frac{Q}{4 \pi \varepsilon_0 R^3}(-\frac{x}{R^3} \mid_R^\infty + \frac{1}{2} x^2\mid_r^R)
    =\frac{Q}{4 \pi \varepsilon_0 R ^3} \frac{3R^2-r^2}{2}$
    
    $ E_\phi =\frac{1}{2} \int \rho u d\tau$

    $= \frac{1}{4} \frac{Q}{4 \pi \varepsilon_0 R^3} \frac{3Q}{4 \pi R^3} \int_{\Omega} (3R^2-r^2)dV$

   \end{spacing}
   \vspace{10pt}
    $=\frac{3 Q^2}{20 \pi \varepsilon_0 R }$
    

\end{enumerate}
\subsubsection{电场的能量}
在某一带电球壳周围存在伴存电场，假设其等势面是一系列封闭曲面，空间被这些等势面划分成许多薄层，如图所示：

将带点球壳每一部分从无穷远处移动到原处，在移动dx段距离的过程中，只有这一部分的电场发生了变化，故外力所做的功存储在这一部分空间中
两等势面的电势差为$dU$,每一薄层的电能$dW=\frac{1}{2} QdU$,整个电场的能量$W=\frac{1}{2}Q \int dU=\frac{1}{2}QU_0$

根据高斯定理有：$\oint \vec{E}dS=\frac{1}{\varepsilon_0} Q$

设$dl$和$dS$的方向一致（即指向S的外法线方向）则有：$dU=\vec{E} dl$

所以：$W=\oint \frac{1}{2}\varepsilon_0 \vec{E} \cdot \vec{E}d\tau$

所以单位体积内的电场能量为：
\begin{equation}
w_e=\frac{1}{2} \varepsilon_0 E^2
\end{equation}

$w_e$叫做电场的能量密度,虽然是由孤立带点等势面的特例得到，却对所有电场都适用

例题：

\begin{itemize}
\item 有一个厚度为t的球壳，内半径为R，将无穷远处的一个点电荷移动到导体空腔圆心处，求在移动过程中外力所做的功

解：

方法一：

    从初态到末态，可以看作只有导体所占空间处的电场消失了，用电场的能量密度求解：

$W=\frac{1}{2} \varepsilon_0 \int_R^{R+t}\vec{E}^2 d\tau 
=\frac{Q^2}{32 \pi^2 \varepsilon_0} \int_R^{R+t}\frac{1}{r^4} d\tau 
=\frac{Q^2}{8 \pi \varepsilon_0}(\frac{1}{R+t}-\frac{1}{R})$

方法二：
    求移动后体系的静电能

\end{itemize}
\subsubsection{电子的经典半径}
带电粒子和他的电场是不可分割的整体，若电场的能量为W,按狭义相对论粒子的质量为
\begin{equation}
m_{em}=\frac{W}{c^2}
\end{equation}
其中，$m_{em}$是粒子的电磁质量，$W$是电子的电场能量，$c$是光速。
$W$与电荷的分布有关，

假设粒子半径为$r_0$，电荷密度为$\rho$，电荷量为$q$，

电荷均匀分布在球面上：
\begin{equation}
    m_{em}=\frac{1}{2} \frac{Q^2}{4 \pi \varepsilon_0 c^2 R}
\end{equation}

电荷均匀分布在球体上：
\begin{equation}
    m_{em}=\frac{3}{5} \frac{Q^2}{4 \pi \varepsilon_0 c^2 R}
\end{equation}

迄今还无法测知电荷的分布，作为一种估计，$m_{em}=\frac{Q^2}{4 \pi \epsilon_0 c^2 R}
$
由上式可得：
\begin{equation}
    R=\frac{Q^2}{4 \pi \varepsilon_0 c^2 m_{em}}
\end{equation}
粒子的电量是可以测得的，但是$m_{em}$是未知的

电子的电荷量$q=-1.6 \times 10^{-19} C$，质量$m_e=9.1 \times 10^{-31} kg$，
若将其全部估算为电磁质量，则
\begin{equation}
    R=\frac{(1.6 \times 10^{-19})^2}{4 \pi \varepsilon_0 c^2 m_e} \approx 2.82 \times 10^{-13} m
\end{equation}
即电子的经典半径，是根据宏观电磁学对电子半径的估计，不可能是准确的。

\subsection{导体的静电平衡和电子的发射}
\subsubsection{导体中的自由电子}
势阱： 在导体中，电荷分布是均匀的，电场强度为0，电势是常数，

逸出功：

\subsubsection{静电平衡}
均匀导体板在匀强电场中的情况

不规则导体在电场中的情况

导体有空腔的情况：空腔内部电场为0

初始外电场为0，导体有空腔，空腔内有电荷：

同上，但导体接地：

同上，在导体外部有电荷：

\section{电流场}
\subsection{导体中的传导电流}
\subsection{电源及其电动势}
\subsection{电容}
\subsection{电流场的连续性}


\end{document}